\begin{table}[tbp]
\centering
\small
\caption{All category-control contrasts and descriptive crossed intervals.}
\label{tab:categorygains}
\begin{tabular}{rllrrl}
\toprule
R & Titles & EEG & Category MAE & EEG MAE & Gain [95\% interval] \\
\midrule
0 & Unfamiliar & Spectral & 1.222 & 1.259 & -0.037 [-0.244, +0.153] \\
0 & Unfamiliar & EEGNet & 1.222 & 1.191 & +0.031 [-0.194, +0.285] \\
0 & Familiar & Spectral & 1.162 & 1.198 & -0.036 [-0.208, +0.172] \\
0 & Familiar & EEGNet & 1.162 & 1.259 & -0.097 [-0.279, +0.076] \\
1 & Unfamiliar & Spectral & 1.232 & 1.177 & +0.055 [-0.166, +0.271] \\
1 & Unfamiliar & EEGNet & 1.232 & 1.311 & -0.079 [-0.298, +0.141] \\
1 & Familiar & Spectral & 1.220 & 1.216 & +0.005 [-0.189, +0.184] \\
1 & Familiar & EEGNet & 1.220 & 1.297 & -0.077 [-0.292, +0.110] \\
2 & Unfamiliar & Spectral & 1.327 & 1.211 & +0.116 [-0.079, +0.327] \\
2 & Unfamiliar & EEGNet & 1.327 & 1.263 & +0.064 [-0.091, +0.232] \\
2 & Familiar & Spectral & 1.134 & 1.214 & -0.080 [-0.256, +0.097] \\
2 & Familiar & EEGNet & 1.134 & 1.320 & -0.185 [-0.350, -0.014] \\
\bottomrule
\end{tabular}
\par\medskip\begin{minipage}{0.97\linewidth}\footnotesize Gain is category-control MAE minus EEG-recipe MAE. It compares whole prediction recipes with different metadata/readout forms. Intervals use 4,000 crossed person/title resamples (seed 425401), reused development and no selection correction. An interval including zero does not establish equivalence or lack of EEG information.\end{minipage}
\end{table}
